\section{Limitations and Next Experiments}
CRR is intended as a reliability framework, not as a universal accuracy replacement. First, the method is not an MVTec AUROC SOTA claim. Official SubspaceAD remains a very strong ranking baseline, and our MVTec contribution is mainly storage, calibration, and diagnostics. Second, CRR is not adversarially robust: label-aware PCA-residual FGSM tests show severe collapse and objective sensitivity. Third, conformal reliability depends on whether the few normal support images are representative of future normal data. If support normals are not exchangeable with test normals, p-values can become overconfident. Fourth, SC3R's evidence has boundaries: the gate-passing result covers all 15 MVTec classes but only $k{=}4$ and seeds 0--2, per-condition false-alarm rates slightly exceed the $\alpha{+}0.02$ budget for JPEG at $\alpha{=}0.10$ (0.121) and for Gaussian noise/JPEG at $\alpha{=}0.20$ (0.226/0.235), it requires a multi-category deployment with matched acquisition conditions, and cross-dataset transfer of source archives is untested. Fifth, ECE on conformal-derived scores is strongly prevalence-sensitive, so calibration numbers on anomaly-balanced test sets must not be read as deployment reliability; we treat operational false-alarm, power, and precision metrics as primary.

The remaining extensions are scoped: SC3R at $k{=}8$ and with more seeds, cross-dataset source archives (e.g., MVTec sources for VisA targets), randomized/smoothed few-shot p-values as an alternative route below the attainable-alpha floor, temperature scaling, isotonic regression, and histogram binning as additional label-free calibration baselines under the same synthetic protocol, and a vision--language few-shot baseline such as WinCLIP~\cite{winclip2023} run under our audit rather than cited from reported numbers.
